\sgaStableTarget{sga6.e02.expose}\section*{第 II 讲. 整体消解的存在性}
\begin{center}
L. Illusie
\end{center}

\sgaStableTarget{sga6.e02.sec.1}\section*{1. 一般整体化判据}

\sgaStableTarget{sga6.e02.sec.1.0}\subsection*{1.0}
本节给定一个景（site）$S$, 一个平坦阿贝尔 $S$-范畴 $\mathcal C$（\hyperlink{sga6.e01.sec.1.1}{I 1.1}）和一个严格满的加性子 $S$-范畴 $\mathcal C_0$（\hyperlink{sga6.e01.sec.1.1}{I 1.1}）. 假设 $\mathcal C_0$ 在 $\mathcal C$ 中局部拟可提升（\hyperlink{sga6.e01.sec.1.2}{I 1.2}）, 并且对满射的核局部拟稳定（\hyperlink{sga6.e01.sec.1.2}{I 1.2}）.

\sgaStableTarget{sga6.e02.prop.1-1}\textbf{命题 1.1.} 设 $S\in \Ob(S)$, 并假设:
\begin{enumerate}[label=(\roman*)]
\item $\mathcal C_{0S}$ 在 $\mathcal C_S$ 中拟可提升（\hyperlink{sga6.e01.sec.1.2}{I 1.2}）, 其中将 $\mathcal C_S$ 看作单点景上的纤维范畴;
\item $\mathcal C_S$ 中每个有限 $\mathcal C_0$-型（\hyperlink{sga6.e01.sec.1.2}{I 1.2}）, 分别为有限 $\mathcal C_0$-表现（\hyperlink{sga6.e01.def.2-8}{I 2.8}）的对象, 都是有限 $\mathcal C_{0S}$-型的.
\end{enumerate}
则:

\begin{enumerate}[label=\alph*)]
\item $\mathcal C_{0S}$ 对满射的核拟稳定（\hyperlink{sga6.e01.sec.1.2}{I 1.2}）, 其中将 $\mathcal C_{0S}$ 看作单点景上的纤维范畴.
\item 设 $E\in\Ob \mathrm D^{-}(\mathcal C_S)$ 且 $n\in\mathbb Z$. 若 $E$ 是 $n$-$\mathcal C_0$-伪凝聚的（\hyperlink{sga6.e01.prop.2-2}{I 2.2}）, 则 $E$ 是 $n$-$\mathcal C_{0S}$-伪凝聚的, 分别为 $(n+1)$-$\mathcal C_{0S}$-伪凝聚的. 换言之,
\[
\mathrm D^{-}(\mathcal C_S)_{n-\mathcal C_0\text{-coh}}
=
\mathrm D^{-}(\mathcal C_S)_{n-\mathcal C_{0S}\text{-coh}},
\]
以及, 分别有
\[
\mathrm D^{-}(\mathcal C_S)_{n-\mathcal C_0\text{-coh}}
\subset
\mathrm D^{-}(\mathcal C_S)_{(n+1)-\mathcal C_{0S}\text{-coh}}
\subset
\mathrm D^{-}(\mathcal C_S)_{(n+1)-\mathcal C_0\text{-coh}},
\]
这些是 $\mathrm D(\mathcal C_S)$ 的严格满子范畴之间的包含关系, 右边的包含关系是平凡的.
\item 对象 $E\in\Ob\mathrm D(\mathcal C_S)$ 关于 $\mathcal C_0$ 伪凝聚（\hyperlink{sga6.e01.def.2-3}{I 2.3}）的充要条件是它关于 $\mathcal C_{0S}$ 伪凝聚; 等价地,
\[
\mathrm D^{-}(\mathcal C_S)_{\mathcal C_0\text{-coh}}
=
\mathrm D^{-}(\mathcal C_S)_{\mathcal C_{0S}\text{-coh}}.
\]
此外, 典范函子
\[
\mathrm K^{-}(\mathcal C_{0S})/\mathrm K^{-,\delta}(\mathcal C_{0S})
\longrightarrow \mathrm D(\mathcal C_S),
\]
是全忠实的, 且其本质像为 $\mathrm D^{-}(\mathcal C_S)_{\mathcal C_0\text{-coh}}$; 这里 $\mathrm K^{-,\delta}(\mathcal C_{0S})$ 表示 $\mathrm K^{-}(\mathcal C_{0S})$ 中由无上同调复形组成的厚子范畴. 若 $\mathcal C_{0S}$ 在 $\mathcal C_S$ 中可提升（\hyperlink{sga6.e01.sec.1.2}{I 1.2}）, 即 $\mathcal C_{0S}$ 的对象都是投射的, 则 $\mathrm K^{-,\delta}(\mathcal C_{0S})$ 仅由零对象组成, 因而函子
\[
\mathrm K^{-}(\mathcal C_{0S})\longrightarrow \mathrm D(\mathcal C_S)
\]
是全忠实的, 且其本质像为 $\mathrm D^{-}(\mathcal C_S)_{\mathcal C_0\text{-coh}}$.
\end{enumerate}

\emph{证明.} a) 设
\[
0\longrightarrow L\longrightarrow L^0\longrightarrow\cdots\longrightarrow L^m\longrightarrow0
\]
是 $\mathcal C_S$ 中的正合列, 对所有 $i$ 都有 $L^i\in\Ob\mathcal C_{0S}$. 只须证明 $L$ 是有限 $\mathcal C_{0S}$-型的. 由（\hyperlink{sga6.e01.prop.2-2}{I 2.2}）, $L$ 是 $\mathcal C_0$-伪凝聚的, 因而由（\hyperlink{sga6.e01.prop.2-9}{I 2.9}）它是有限 $\mathcal C_0$-表现的; 断言于是由 (ii) 得出.

b) 对 $i$ 作递降归纳, 证明当 $i\ge n$, 分别当 $i\ge n+1$ 时, $E$ 是 $i$-$\mathcal C_{0S}$-伪凝聚的. 可以先假设 $E$ 是严格 $(n+1)$-$\mathcal C_{0S}$-伪凝聚的, 分别为严格 $(n+2)$-$\mathcal C_{0S}$-伪凝聚的（\hyperlink{sga6.e01.def.2-1}{I 2.1}）; 于是只须证明 $E$ 是 $n$-$\mathcal C_{0S}$-伪凝聚的, 分别为 $(n+1)$-$\mathcal C_{0S}$-伪凝聚的. 因而有正合列
\[
0\longrightarrow E'\longrightarrow E\longrightarrow E''\longrightarrow0,
\]
其中 $E'$ 是严格 $\mathcal C_{0S}$-完美的（\hyperlink{sga6.e01.def.2-1}{I 2.1}）, 且当 $i>n$, 分别当 $i>n+1$ 时, $E''{}^i=0$. 因此以 $E''$ 代替 $E$, 可以假设（\hyperlink{sga6.e01.schol.2-6}{I 2.6}）当 $i>n$, 分别当 $i>n+1$ 时有 $E^i=0$. 于是由（\hyperlink{sga6.e01.prop.2-10}{I 2.10} a)）, $H^n(E)$, 分别为 $H^{n+1}(E)$, 是有限 $\mathcal C_0$-型的, 分别为有限 $\mathcal C_0$-表现的. 因而由 (ii), $H^n(E)$, 分别为 $H^{n+1}(E)$, 是有限 $\mathcal C_{0S}$-型的. 再用（\hyperlink{sga6.e01.prop.2-10}{I 2.10} b)）, 得到 $E$ 是 $n$-$\mathcal C_{0S}$-伪凝聚的, 分别为 $(n+1)$-$\mathcal C_{0S}$-伪凝聚的. 在引用（\hyperlink{sga6.e01.schol.2-6}{I 2.6} 和 \hyperlink{sga6.e01.prop.2-10}{2.10}）时, 已隐含使用假设 (i) 和 a) 的结论.

c) 第一条断言由 b) 得出, 第二条由（\hyperlink{sga6.e01.prop.2-7}{I 2.7}）得出, 最后一条是显然的.

\sgaStableTarget{sga6.e02.prop.1-2}\textbf{命题 1.2.} 假设 $\mathcal C_0$ 在 $\mathcal C$ 中局部可提升（\hyperlink{sga6.e01.sec.1.2}{I 1.2}）, 且对满射的核局部稳定（\hyperlink{sga6.e01.sec.1.2}{I 1.2}）, 并作出（\hyperlink{sga6.e02.prop.1-1}{1.1 (i)} 以及分别的 (ii)）中的假设. 以 $\mathcal C_{1S}$ 表示 $\mathcal C_S$ 中由局部属于 $\mathcal C_0$ 的对象组成的严格满子范畴. 则:
\begin{enumerate}[label=\alph*)]
\item $\mathcal C_{1S}$ 对满射的核稳定（\hyperlink{sga6.e01.sec.1.2}{I 1.2}）, 且作为单点景上的纤维范畴, 在 $\mathcal C_S$ 中拟可提升（\hyperlink{sga6.e01.sec.1.2}{I 1.2}）.
\item 设 $E\in\Ob\mathrm D(\mathcal C_S)$, 并设 $[m,n]$ 是 $\mathbb Z$ 的区间. $E$ 的 $\mathcal C_0$-完美幅包含在 $[m,n]$ 中的充要条件是: $E$ 在 $\mathrm D(\mathcal C_S)$ 中同构于一个复形 $F$, 满足当 $i\notin[m,n]$ 时 $F^i=0$, 当 $i>m$ 时 $F^i\in\Ob\mathcal C_{0S}$, 且 $F^m\in\Ob\mathcal C_{1S}$.
\item 典范函子
\[
\mathrm K^b(\mathcal C_{1S})/\mathrm K^{b,\delta}(\mathcal C_{1S})
\longrightarrow \mathrm D(\mathcal C_S)
\]
是全忠实的, 且其本质像为 $\mathrm D(\mathcal C_S)$ 中由有限 $\mathcal C_0$-完美幅的复形组成的满子范畴.
\item 若 $\mathcal C_{1S}$ 在 $\mathcal C_S$ 中不只是拟可提升, 而且可提升, 则 b) 中的“$E$ 在 $\mathrm D(\mathcal C_S)$ 中同构于 $F$”可换成“存在拟同构 $F\to E$”; 并且 c) 中的 $\mathrm K^{b,\delta}(\mathcal C_{1S})$ 仅由零对象组成.
\end{enumerate}

\emph{证明.} a) 显然 $\mathcal C_{1S}$ 对满射的核稳定. 另一方面, $\mathcal C_{1S}$ 的对象是 $\mathcal C_0$-伪凝聚的, 因而由（\hyperlink{sga6.e02.prop.1-1}{1.1} c)）也是 $\mathcal C_{0S}$-伪凝聚的; 取 $S$ 为单点景, 第二条断言便由（\hyperlink{sga6.e01.cor.2-14}{I 2.14}）得出.

b) 只须证明必要性. 由（\hyperlink{sga6.e02.prop.1-1}{1.1} c)）, $E$ 已经是 $\mathcal C_{0S}$-伪凝聚的. 由于 $E$ 在次数 $>n$ 时无上同调, 在以 $\mathrm D(\mathcal C_S)$ 中同构的复形代替 $E$ 后, 可以假设（\hyperlink{sga6.e01.prop.2-10}{I 2.10}）$E$ 是严格 $(m+1)$-$\mathcal C_{0S}$-伪凝聚的, 且当 $i>n$ 时有 $E^i=0$. 由于 $E$ 在次数 $<m$ 时无上同调, 还可以用其截断
\[
0\longrightarrow E^m/B^m\longrightarrow E^{m+1}\longrightarrow\cdots\longrightarrow E^n\longrightarrow0,
\]
代替 $E$, 使得当 $i<m$ 时有 $E^i=0$. 于是由（\hyperlink{sga6.e01.lem.4-16}{I 4.16}）, $E^m$ 局部是 $\mathcal C_0$ 的对象, 因而是 $\mathcal C_{1S}$ 的对象.

c) 由（\hyperlink{sga6.e02.prop.1-2}{1.2} a)）, 断言可由（\hyperlink{sga6.e02.prop.1-2}{1.2} b)）及（\hyperlink{sga6.e01.rem.4-12-2}{I 4.12.2}）得出. d) $\mathrm K^{b,\delta}(\mathcal C_{1S})$ 仅由零对象组成是显然的. 另一断言由应用于单点景的（\hyperlink{sga6.e01.cor.4-6}{I 4.6}）得出.

\sgaStableTarget{sga6.e02.sec.2}\section*{2. 对若干模范畴的应用}

\sgaStableTarget{sga6.e02.sec.2.0}\subsection*{2.0. 记号}
设 $(S,A)$ 是赋环拓扑斯. 对 $U\in\Ob S$, 以 $\Mod(U)$ 表示左 $A_U$-模的范畴; 以 $\Lib(U)$, 分别以 $\Loclib(U)$, 表示 $\Mod(U)$ 中由有限型自由 $A_U$-模, 分别由局部为有限型自由模的直和因子的模组成的满子范畴.

\sgaStableTarget{sga6.e02.sec.2.1}\subsection*{2.1. 提升引理}

\sgaStableTarget{sga6.e02.lem.2-1-1}\textbf{引理 2.1.1}（参见 EGA $0_{\mathrm I}$, 5.2.3）. 设 $(S,A)$ 是赋环拓扑斯, 且 $S\in\Ob S$. 假设 $S$ 拟紧, 即每个覆盖族 $(S_i\to S)$ 都由某个有限覆盖族控制.
\begin{enumerate}[label=\alph*)]
\item 设左 $A_S$-模 $F$ 是以滤过范畴 $\Lambda$ 为指标的 $A_S$-模 $F_\lambda$ 的滤过归纳极限, 并设
\[
u=\varinjlim u_\lambda:F\longrightarrow G
\]
为满射, 其中 $G$ 是有限型 $A_S$-模. 则存在 $\lambda\in\Lambda$ 使 $u_\lambda$ 为满射.
\item 设 $F$ 是由其截面生成的左 $A_S$-模, 即对适当的指标集 $I$ 存在满射 $A_S^{(I)}\to F$. 若 $F$ 是有限型的, 则对某个 $p\in\mathbb N$ 存在满射 $A_S^p\to F$.
\end{enumerate}
a) 的证明类似于所引结果的证明, 留给读者; b) 立即由 a) 得出.

\sgaStableTarget{sga6.e02.lem.2-1-2}\textbf{引理 2.1.2.} 设 $(S,A)$ 是赋环拓扑斯, $\mathcal C$ 是以左 $A$-模为纤维的 $S$-范畴, $\mathcal C_0$ 是由有限型自由 $A$-模组成的纤维子范畴（\hyperlink{sga6.e01.sec.1.3}{I 1.3} d)）, $S\in\Ob S$, 而 $\mathcal C'_S$ 是 $\mathcal C_S$ 的严格满加性子范畴, 包含 $\mathcal C_{0S}$, 且对核、余核及扩张稳定. 假设
\[
H^1(S,F)=0\quad\text{对每个 }F\in\Ob\mathcal C'_S.
\]
则:
\begin{enumerate}[label=\alph*)]
\item $\mathcal C_S$ 中每个满足下列条件的正合列
\[
0\longrightarrow E\longrightarrow F\longrightarrow G\longrightarrow0,
\]
都分裂: $\operatorname{Hom}(G,E)\in\Ob\mathcal C'_S$, 且 $G$ 局部是有限型自由模的直和因子. 范畴 $\mathcal C_{0S}$ 在 $\mathcal C'_S$ 中可提升（\hyperlink{sga6.e01.prop.2-2}{I 2.2}）.
\item 以 $\Mod(A(S))$ 表示 $A(S)$-模的范畴. 函子
\[
\Gamma_S:\mathcal C'_S\longrightarrow\Mod(A(S))
\]
是正合的. 此外:
\begin{enumerate}[label=(\roman*)]
\item 若 $\mathcal C'_S$ 的每个对象 $F$ 都有有限整体表现
\[
A_S^q\longrightarrow A_S^p\longrightarrow F\longrightarrow0,
\]
则 $\Gamma_S$ 是全忠实的, 其本质像为 $\Mod(A(S))$ 中由有限表现 $A(S)$-模组成的满子范畴;
\item 若 $\mathcal C'_S$ 包含所有自由 $A_S$-模, 若 $\mathcal C'_S$ 的每个对象 $F$ 都有整体表现
\[
A_S^{(J)}\longrightarrow A_S^{(I)}\longrightarrow F\longrightarrow0,
\]
且 $\Gamma_S$ 与滤过归纳极限交换, 则 $\Gamma_S$ 是范畴等价.
\end{enumerate}
\end{enumerate}

\emph{证明.} a) 由于 $G$ 局部是有限型自由模的直和因子, 谱序列
\[
E_2^{pq}=H^p(S,\mathcal{E}xt^q(G,E))\Longrightarrow \Ext^{p+q}(G,E)
\]
退化（\hyperlink{sga6.e01.lem.1-3-1}{I 1.3.1}）, 并给出同构
\[
H^1(S,\operatorname{Hom}(G,E))\simeq \Ext^1(G,E),
\]
由此得到第一条断言. 第二条是其特例, 因为当 $G\in\Ob\mathcal C_{0S}$ 且 $E\in\Ob\mathcal C'_S$ 时, 有 $\operatorname{Hom}(G,E)\in\Ob\mathcal C'_S$.

b) 第一条断言是平凡的. 下面在情形 (i) 或 (ii) 中证明 $\Gamma_S:\mathcal C'_S\to\Mod(A(S))$ 的全忠实性. 设 $F,G\in\Ob\mathcal C'_S$, 并设
\sgaStableTarget{sga6.e02.eq.unnamed}\[
A_S^{(J)}\longrightarrow A_S^{(I)}\longrightarrow F\longrightarrow0
\tag{$*$}
\]
是 $F$ 的一个整体表现, 在情形 (i) 中 $I$ 和 $J$ 有限. 应用 $\Gamma_S$, 由于它在 $\mathcal C'_S$ 上正合, 且在情形 (ii) 中与滤过归纳极限交换, 得到正合列
\sgaStableTarget{sga6.e02.eq.unnamed.occ2}\[
A(S)^{(J)}\longrightarrow A(S)^{(I)}\longrightarrow\Gamma_S(F)\longrightarrow0.
\tag{$**$}
\]
将左正合函子 $\operatorname{Hom}_{\mathcal C_S}(-,G)$ 和 $\operatorname{Hom}_{A(S)}(-,\Gamma_S(G))$ 分别应用于 $(*)$ 和 $(**)$, 得到行正合的交换图:
\[
\begin{tikzcd}[column sep=small]
0 \arrow[r] & \operatorname{Hom}_{\mathcal C_S}(F,G) \arrow[r] \arrow[d] &
\operatorname{Hom}_{\mathcal C_S}(A_S^{(I)},G) \arrow[r] \arrow[d,"\sim"] &
\operatorname{Hom}_{\mathcal C_S}(A_S^{(J)},G) \arrow[d,"\sim"] \\
0 \arrow[r] & \operatorname{Hom}_{A(S)}(\Gamma_SF,\Gamma_SG) \arrow[r] &
\operatorname{Hom}_{A(S)}(A(S)^{(I)},\Gamma_SG) \arrow[r] &
\operatorname{Hom}_{A(S)}(A(S)^{(J)},\Gamma_SG).
\end{tikzcd}
\]
右边两个竖直箭头显然是同构, 因而左边的箭头也是同构. 情形 (i) 和 (ii) 中本质像的描述是直接的, 留给读者.

\sgaStableTarget{sga6.e02.lem.2-1-3}\textbf{引理 2.1.3.} 设 $\mathcal A$ 和 $\mathcal B$ 是阿贝尔范畴, $\varphi:\mathcal A\to\mathcal B$ 是正合函子.
\begin{enumerate}[label=\alph*)]
\item 对 $M\in\Ob\mathcal A$, 下列条件等价:
\begin{enumerate}[label=(\roman*)]
\item 对每个 $L\in\Ob\mathcal A$, 由 $\varphi$ 诱导的典范映射
\[
\Ext^n_{\mathcal A}(L,M)\longrightarrow\Ext^n_{\mathcal B}(\varphi L,\varphi M)
\]
对所有 $n$ 都是同构;
\item 对每个 $L\in\Ob\mathcal A$, 典范映射
\[
\operatorname{Hom}_{\mathcal A}(L,M)\longrightarrow\operatorname{Hom}_{\mathcal B}(\varphi L,\varphi M)
\]
是同构, 且对每个 $n\ge1$ 及每个 $x\in\Ext^n_{\mathcal B}(\varphi L,\varphi M)$, 在 $\mathcal A$ 中都存在一个消去 $x$ 的满射 $L'\to L$.
\end{enumerate}
\item 下列条件等价:
\begin{enumerate}[label=(\roman*)]
\item 对所有 $L,M\in\Ob\mathcal A$ 及所有 $n\in\mathbb Z$, 典范映射
\[
\Ext^n_{\mathcal A}(L,M)\longrightarrow\Ext^n_{\mathcal B}(\varphi L,\varphi M)
\]
是同构;
\item 由 $\varphi$ 诱导的典范函子
\[
\mathrm D^b(\mathcal A)\longrightarrow\mathrm D^b(\mathcal B)
\]
是全忠实的, 且其本质像为 $\mathrm D^b(\mathcal B)$ 中由满足如下条件的复形 $F$ 组成的满子范畴: 对每个 $i$, $H^i(F)$ 都属于 $\varphi$ 的本质像.
\end{enumerate}
\end{enumerate}

\emph{证明.} 对 a) 中的 (i)$\Rightarrow$(ii), 只须证明: 若 $n\ge1$, $L\in\Ob\mathcal A$, 且 $y\in\Ext^n_{\mathcal A}(L,M)$, 则存在满射 $L'\to L$ 消去 $y$. 利用 Yoneda 的解释（[18], III 3.2）, $y$ 是一个正合列
\[
0\longrightarrow M\longrightarrow E_{n-1}\longrightarrow\cdots\longrightarrow E_0\longrightarrow L\longrightarrow0,
\]
的类, 而满射 $E_0\to L$ 消去 $y$（[18], III 3.2.5 及 3.2.8）. 反之, 令
\[
F^n(L)=\Ext^n_{\mathcal A}(L,M),\qquad
G^n(L)=\Ext^n_{\mathcal B}(\varphi L,\varphi M).
\]
可消去性假设表明, 对 $n\ge1$ 及 $L\in\Ob\mathcal A$, 有
\[
G^n(L)=\varinjlim\operatorname{Coker}\bigl(G^{n-1}(L_0)\to G^{n-1}(L_1)\bigr),
\]
其中极限取遍正合列 $L_1\to L_0\to L\to0$. 如上所述, $F^n(L)$ 也满足同样的公式. 由于 $F^0\to G^0$ 是同构, 对 $n$ 归纳即知由 $\varphi$ 诱导的态射 $F^n\to G^n$ 对每个 $n$ 都是同构. b) 中的等价性由直接的逐步分解论证得出, 留给读者.

\sgaStableTarget{sga6.e02.rem.2-1-3-1}\textbf{注 2.1.3.1.} a) 当然有一个对偶版本, 其表述留给读者. 另一方面, 若 $\mathcal A$ 和 $\mathcal B$ 有足够多的内射对象, 则 b) 中的条件 (i)、(ii) 等价于:
\begin{enumerate}[label=(\roman*),start=3]
\item 对 $L\in\Ob\mathrm D^{-}(\mathcal A)$ 及 $M\in\Ob\mathrm D^{+}(\mathcal A)$, 典范映射
\[
\mathrm R\!\operatorname{Hom}(L,M)\longrightarrow
\mathrm R\!\operatorname{Hom}(\varphi L,\varphi M)
\]
是同构,
\end{enumerate}
这由谱序列论证立即可见.

\sgaStableTarget{sga6.e02.sec.2.2}\subsection*{2.2. 诺特概形与除子性概形}

\sgaStableTarget{sga6.e02.notation.2-2-1}\textbf{记号 2.2.1.} 设 $S$ 是概形. 以 $\Qcoh(S)$, 分别以 $\Coh(S)$, 表示拟凝聚, 分别凝聚的 $\mathcal O_S$-模所成的阿贝尔范畴. 设 $[m,n]$ 是 $\mathbb Z$ 的区间. 若 $\mathrm D(\Qcoh(S))$ 的对象 $E$ 关于 $S$ 的 Zariski 景上的纤维范畴
\[
U\longmapsto \Loclib(U)
\]
是 $n$-伪凝聚的, 分别为伪凝聚的、完美幅包含在 $[m,n]$ 中的、完美的, 则相应地称 $E$ 为 $n$-伪凝聚的、伪凝聚的、完美幅包含在 $[m,n]$ 中的、完美的（（\hyperlink{sga6.e01.def.2-3}{I 2.3}）、（\hyperlink{sga6.e01.def.4-7}{I 4.7}）及 \hyperlink{sga6.e02.sec.2.0}{(2.0)}）. 以 $\mathrm D^{-}(\Qcoh(S))_{\coh}$, 分别以 $\mathrm D(\Qcoh(S))_{\parf}$, 表示 $\mathrm D(\Qcoh(S))$ 中由伪凝聚, 分别由完美复形组成的满三角子范畴.

\sgaStableTarget{sga6.e02.prop.2-2-2}\textbf{命题 2.2.2.} 设 $S$ 是诺特概形. 典范函子
\[
\mathrm D^{-}(\Coh(S))\longrightarrow \mathrm D(\Qcoh(S))
\]
是全忠实的, 且其本质像为 $\mathrm D^{-}(\Qcoh(S))_{\coh}$.

\emph{证明.} 应用 \hyperlink{sga6.e02.prop.1-1}{(1.1)}, 取景为 $S$ 的 Zariski 景, 取 $\mathcal C$, 分别取 $\mathcal C_0$, 为纤维范畴 $U\mapsto\Qcoh(U)$, 分别为 $U\mapsto\Coh(U)$. 由于 $\mathcal O_S$ 凝聚, $\mathcal C_0$ 的对象伪凝聚（\hyperlink{sga6.e01.cor.3-5}{I 3.5}）; 因而 $\mathrm D(\Qcoh(U))$ 的对象在 \hyperlink{sga6.e02.notation.2-2-1}{(2.2.1)} 的意义下伪凝聚, 当且仅当它关于 $\mathcal C_0$ 伪凝聚（\hyperlink{sga6.e01.schol.2-14-1}{I 2.14.1}）. 由（\hyperlink{sga6.e02.prop.1-1}{1.1} c)）, 只须验证 \hyperlink{sga6.e02.prop.1-1}{(1.1)} 的条件 (i)、(ii). 对 (ii), 这是平凡的, 因为凝聚模的拟凝聚商模凝聚, 且凝聚性是局部性质. 对 (i), 使用（\hyperlink{sga6.e02.lem.2-1-1}{2.1.1} a)）及下述事实: 每个拟凝聚 $\mathcal O_S$-模都是凝聚 $\mathcal O_S$-模的滤过归纳极限（EGA I 9.4.9）.

暂且承认与本节无关的 \hyperlink{sga6.e02.prop.3-5}{(3.5)} 和 \hyperlink{sga6.e02.prop.3-7}{(3.7)}, 便容易推出下述结果.

\sgaStableTarget{sga6.e02.cor.2-2-2-1}\textbf{推论 2.2.2.1.} 设 $S$ 是诺特概形, 分别为有限维诺特概形. 典范函子
\[
\mathrm D^b(\Coh(S))\longrightarrow \mathrm D(S)
\]
分别为
\[
\mathrm D^{-}(\Coh(S))\longrightarrow \mathrm D(S)
\]
是全忠实的, 其本质像为 $\mathrm D^b(S)_{\coh}$, 分别为 $\mathrm D^{-}(S)_{\coh}$（\hyperlink{sga6.e01.ex.2-15}{I 2.15}）.

\sgaStableTarget{sga6.e02.prop.2-2-3}\textbf{命题 2.2.3}（参见 EGA II 4.5.2 及 4.5.5）. 设 $S$ 是拟紧拟分离概形, $(L_i)_{i\in I}$ 是 $S$ 上的可逆层族. 对每个 $\mathcal O_S$-模 $F$ 及每对 $(i,n)\in I\times\mathbb Z$, 以 $F_{(i,n)}$ 表示 $F\otimes L_i^{\otimes n}$ 中由它在 $S$ 上的截面生成的子 $\mathcal O_S$-模. 下列条件等价:
\begin{enumerate}[label=(\roman*)]
\item[(i)] 当 $f\in\Gamma(S,L_i^{\otimes n})$, $i\in I$ 且 $n\in\mathbb N$ 时, 开集 $S_f$ 构成 $S$ 的拓扑的一组基;
\item[(i')] 这些 $S_f$ 中为仿射的那些覆盖 $S$;
\item[(ii)] 对每个有限型拟凝聚 $\mathcal O_S$-模 $F$, 都存在非负整数族 $(n_i)_{i\in I}$ 和 $(r_i)_{i\in I}$ 以及满射
\[
\bigoplus_{i\in I}\mathcal O_S^{r_i}\otimes L_i^{\otimes(-n_i)}\longrightarrow F;
\]
\item[(iii)] 对每个 $F\in\Ob\Qcoh(S)$, 有
\[
F=\sum_{(i,n)\in I\times\mathbb N}F_{(i,n)}\otimes L_i^{\otimes(-n)},
\]
这里的和是 $F$ 的子层之和, 未必为直和;
\item[(iii')] 与 (iii) 相同的条件, 但只要求 $F$ 为 $\mathcal O_S$ 的拟凝聚理想.
\end{enumerate}

\emph{证明.} 这是对所引段落的简单改述; 为方便读者, 下面简述其论证. 对 (i)$\Rightarrow$(i'), 给定 $x\in S$, 选取包含 $x$ 的仿射开集 $U$. 若开集 $S_f$ 满足 $x\in S_f\subset U$, 其中 $f\in\Gamma(S,L_i^{\otimes n})$ 且 $n\ge0$, 则由下述引理它必为仿射的.

\sgaStableTarget{sga6.e02.lem.2-2-3-1}\textbf{引理 2.2.3.1}（EGA II 5.5.8）. 设 $X$ 是概形, $L$ 是可逆 $\mathcal O_X$-模, 且 $f\in\Gamma(X,L)$. 典范浸入 $X_f\to X$ 是仿射的.

\emph{证明.} 这是平凡的: 问题在 $X$ 上局部, 因而可以假设 $X$ 仿射且 $L\simeq\mathcal O_X$.

对 (i')$\Rightarrow$(ii), 选取 $S$ 的有限仿射开覆盖 $S_{f_{ij}}$, 其中 $f_{ij}\in\Gamma(S,L_i^{\otimes m_{ij}})$, 并令 $S_i=\bigcup_j S_{f_{ij}}$. 以适当的幂代替 $f_{ij}$ 后, 可以假设对固定的 $i$, 所有 $m_{ij}$ 都等于同一个整数 $m_i$. 层 $F|_{S_{f_{ij}}}$ 由 $S_{f_{ij}}$ 上有限个截面 $h_{ijk}$ 生成. 存在整数 $q_i\ge0$（EGA I 9.3.1 及 IV 1.7.5）, 使得 $h_{ijk}\otimes f_{ij}^{\otimes q_i}$ 可延拓为 $F\otimes L_i^{\otimes q_im_i}$ 在 $S$ 上的截面. 换言之, $F\otimes L_i^{\otimes q_im_i}$ 在 $S_i$ 上由它在 $S$ 上的有限个截面生成. 按照 EGA II 4.5.5 的证明方法, 得到某个整数 $n_i$, 可假设 $n_i\ge0$, 使得对所有 $n\ge n_i$, $F\otimes L_i^{\otimes n}$ 在 $S_i$ 上都由它在 $S$ 上的有限个截面生成. 特别地, 存在整数 $r_i\ge0$ 及态射
\[
\mathcal O_S^{r_i}\otimes L_i^{\otimes(-n_i)}\longrightarrow F
\]
它在 $S_i$ 上诱导满射. 由此立即得出 (i')$\Rightarrow$(ii).

对 (ii)$\Rightarrow$(iii), 由于 $S$ 上每个拟凝聚层都是其有限型拟凝聚子层的归纳极限（EGA I 9.4.9 及 IV 1.7.7）, 只须对有限型的 $F$ 证明 (iii). 此时存在态射
\[
u_i:\mathcal O_S^{r_i}\otimes L_i^{\otimes(-n_i)}\longrightarrow F
\]
其和为满射. 但
\[
\begin{aligned}\operatorname{Im}(u_i)&=\operatorname{Im}\bigl(u_i\otimes\operatorname{Id}(L_i^{\otimes n_i})\otimes\operatorname{Id}(L_i^{\otimes(-n_i)})\bigr)\\&=\operatorname{Im}\bigl(u_i\otimes\operatorname{Id}(L_i^{\otimes n_i})\bigr)\otimes L_i^{\otimes(-n_i)},\end{aligned}
\]
而 $\operatorname{Im}\bigl(u_i\otimes\operatorname{Id}(L_i^{\otimes n_i})\bigr)$ 是 $F_{(i,n_i)}$ 的子层, 由此得到断言. (iii)$\Rightarrow$(iii') 是平凡的. 最后, 对 (iii')$\Rightarrow$(i), 设 $x\in S$, $U$ 是包含 $x$ 的开集, $F$ 是定义 $S-U$ 的拟凝聚理想. 由 (iii'), 存在 $(i,n)\in I\times\mathbb N$ 以及 $F\otimes L_i^{\otimes n}$ 在 $S$ 上的截面 $f$, 使得 $f(x)\ne0$. 于是 $x\in S_f\subset U$, 这就证明了 (i).

\sgaStableTarget{sga6.e02.def.2-2-4}\textbf{定义 2.2.4.} 设 $S$ 是拟紧拟分离概形, $(L_i)_{i\in I}$ 是 $S$ 上的可逆层族. 若 \hyperlink{sga6.e02.prop.2-2-3}{(2.2.3)} 中的等价条件成立, 则称 $L_i$ 构成一个\emph{丰沛族（ample family）}. 特别地, 由单个层 $L$ 组成的族丰沛, 当且仅当 $L$ 在通常意义下丰沛（EGA II 4.5.3）.

\sgaStableTarget{sga6.e02.def.2-2-5}\textbf{定义 2.2.5.} 若概形 $S$ 拟紧拟分离, 且存在可逆 $\mathcal O_S$-模的丰沛族 \hyperlink{sga6.e02.def.2-2-4}{(2.2.4)}, 则称 $S$ 是\emph{除子性的（divisorial）}.

\sgaStableTarget{sga6.e02.lem.2-2-6}\textbf{引理 2.2.6.} 设 $X$ 是局部诺特概形. 设 $U$ 是 $X$ 的开集, 且包含映射 $U\to X$ 仿射; 例如 $U$ 仿射而 $X$ 分离. 则 $X-U$ 纯余维为 $0$ 或 $1$; 即在 $X-U$ 的每个极大点 $x$ 处都有 $\dim\mathcal O_{X,x}=0$ 或 $1$. 若 $U$ 包含 $X$ 的极大点, 则其纯余维为 $1$.

\emph{证明.} 令 $Y=X-U$. 取 $Y$ 的一个极大点 $x$, 并作基变换 $X'=\operatorname{Spec}\mathcal O_{X,x}\to X$; 得到笛卡尔图
\[
\begin{tikzcd}
U' \arrow[r] \arrow[d] & U \arrow[d] \\
X' \arrow[r] & X .
\end{tikzcd}
\]
令 $Y'=X'-U'$. 由于 $x$ 是 $Y$ 的极大点, $Y'$ 就是 $X'$ 的闭点. 由于 $X'$ 仿射, $U'$ 是 $X'$ 的仿射开集. 下面证明 $X'$ 的维数 $d$ 为 $0$ 或 $1$. 事实上, 已知（SGA 2 V 5.7）$H^d_{Y'}(X',\mathcal O_{X'})\ne0$. 但若 $d\ge2$, 则
\[
H^d_{Y'}(X',\mathcal O_{X'})=H^{d-1}(U',\mathcal O_{X'}),
\]
而后一群因 $U'$ 仿射且 $d-1\ge1$ 而为零, 矛盾.

\sgaStableTarget{sga6.e02.rem.2-2-6-1}\textbf{注 2.2.6.1.} 若无 $X$ 分离的假设, 对 $X$ 的仿射开集 $U$, \hyperlink{sga6.e02.lem.2-2-6}{(2.2.6)} 未必成立; 原点加倍的仿射平面（EGA I 5.5.11）便是反例.

\sgaStableTarget{sga6.e02.prop.2-2-7}\textbf{命题 2.2.7.} 设 $S$ 是诺特、分离且局部因子性的概形, 即其局部环均为因子环. 则 $S$ 为除子性概形 (2.2.5). 特别地（EGA IV 21.11.1）:

\sgaStableTarget{sga6.e02.cor.2-2-7-1}\textbf{推论 2.2.7.1.} 每个诺特正则分离概形都是除子性概形.

\emph{2.2.7 的证明.} 可以假设 $S$ 整. 取 $S$ 的有限非空仿射开覆盖 $(S_i)_{i\in I}$. 由 (2.2.6), $S-S_i$ 纯余维为 $1$; 因而由 EGA IV 21.7.4, 它是某个正除子的支集. 换言之, 对每个 $i\in I$, 存在可逆层 $L_i$ 及截面 $f_i\in\Gamma(S,L_i)$, 使得 $S_i=S_{f_i}$. 由 \hyperlink{sga6.e02.def.2-2-4}{(2.2.4)}, 利用 \hyperlink{sga6.e02.prop.2-2-3}{(2.2.3 (i'))}, 族 $(L_i)$ 丰沛.

\sgaStableTarget{sga6.e02.rem.2-2-7-2}\textbf{注 2.2.7.2.} 如 \hyperlink{sga6.e02.lem.2-2-6}{(2.2.6)} 中一样, $S$ 分离的假设是本质的. 再取 $S$ 为原点加倍的仿射平面. 回忆 $S$ 是把仿射平面 $\mathbb A^2_k$ 的两个副本沿原点 $O$ 的补集 $U$ 上的恒等映射粘合而成的; 因而有典范投影
\[
p:S\longrightarrow\mathbb A^2_k
\]
它在 $U$ 上为同构, 而 $p^{-1}(O)$ 由两个点 $O_1$ 和 $O_2$ 组成. $S$ 上的可逆层由 $\mathbb A^2_k$ 上的两个可逆层通过它们在 $U$ 上的同构粘合得到. 由于 $\mathbb A^2_k$ 上的每个可逆层均平凡, $S$ 上的可逆层由某个元素 $g\in\Gamma(U,\mathcal O_U^*)$ 给出. 但 $g$ 必为常数, 因为原点在平面中的深度为 $2$, 故限制映射 $\Gamma(\mathbb A^2_k,\mathcal O_{\mathbb A^2})\to\Gamma(U,\mathcal O_U)$ 是同构. 因而 $S$ 上的每个可逆层均平凡. 另一方面, 容易看出 $p$ 诱导同构
\[
\Gamma(\mathbb A^2_k,\mathcal O_{\mathbb A^2})\xrightarrow{\sim}\Gamma(S,\mathcal O_S).
\]
因此 $\mathcal O_S$ 的每个在 $O_1$ 处为零的整体截面, 在 $O_2$ 处也为零, 从而当 $f\in\Gamma(S,\mathcal O_S)$ 时, 开集 $S_f$ 不构成 $S$ 的拓扑的一组基.

\sgaStableTarget{sga6.e02.thm.2-2-8}\textbf{定理 2.2.8.} 设 $S$ 是除子性概形 \hyperlink{sga6.e02.def.2-2-5}{(2.2.5)}, $(L_i)_{i\in I}$ 是可逆 $\mathcal O_S$-模的丰沛族. 以 $L(S)$ 表示 $\Qcoh(S)$ 中由所有 $L_i^{\otimes(-n)}$（$(i,n)\in I\times\mathbb N$）生成的严格满加性子范畴. 设 $E\in\Ob\mathrm D(\Qcoh(S))$.
\begin{enumerate}[label=\alph*)]
\item 范畴 $L(S)$ 和 $\Loclib(S)$ 在 $\Qcoh(S)$ 中拟可提升（\hyperlink{sga6.e01.sec.1.2}{I 1.2}）, 且对满射的核拟稳定（\hyperlink{sga6.e01.sec.1.2}{I 1.2}）.
\item 设 $n\in\mathbb Z$. $E$ 在 \hyperlink{sga6.e02.notation.2-2-1}{(2.2.1)} 的意义下 $n$-伪凝聚, 分别伪凝聚, 的充要条件是它关于 $L(S)$ 或 $\Loclib(S)$ 具有相应性质（\hyperlink{sga6.e01.def.2-3}{I 2.3}）. 典范函子
\[
\mathrm K^{-}(L(S))/\mathrm K^{-,\delta}(L(S))\longrightarrow \mathrm D(\Qcoh(S)),
\]
\[
\mathrm K^{-}(\Loclib(S))/\mathrm K^{-,\delta}(\Loclib(S))\longrightarrow \mathrm D(\Qcoh(S))
\]
都是全忠实的, 且其本质像为 $\mathrm D^{-}(\Qcoh(S))_{\coh}$.
\item 设 $[m,n]$ 是 $\mathbb Z$ 的区间. $E$ 的完美幅包含在 $[m,n]$ 中 \hyperlink{sga6.e02.notation.2-2-1}{(2.2.1)} 的充要条件是: $E$ 在 $\mathrm D(\Qcoh(S))$ 中同构于一个复形 $F$, 满足当 $i\notin[m,n]$ 时 $F^i=0$, 当 $i>m$ 时 $F^i\in\Ob L(S)$, 且 $F^m$ 局部自由且有限型. 典范函子
\[
\mathrm K^b(\Loclib(S))/\mathrm K^{b,\delta}(\Loclib(S))\longrightarrow \mathrm D(\Qcoh(S))
\]
是全忠实的, 且其本质像为 $\mathrm D(\Qcoh(S))_{\parf}$.
\end{enumerate}

\emph{证明.} 设 $F\to G$ 为 $\Qcoh(S)$ 中的满射, 其中 $G$ 有限型. 由于 $F$ 是其有限型拟凝聚子层的归纳极限（EGA I 9.4.9）, 由（\hyperlink{sga6.e02.lem.2-1-1}{2.1.1} a)）, 存在 $F$ 的有限型拟凝聚子层 $F'$, 使得 $F'\to F\to G$ 为满射. 由 \hyperlink{sga6.e02.prop.2-2-3}{(2.2.3 (ii))}, 存在满射 $E\to F'$, 其中 $E\in\Ob L(S)$; 因而 $E\to F'\to F\to G$ 为满射. 这就证明有限型拟凝聚 $\mathcal O_S$-模的范畴在 $\Qcoh(S)$ 中拟可提升, 因而 $\Loclib(S)$ 和 $L(S)$ 更是如此.

对 a) 的第二条断言, 在 $S$ 的 Zariski 景上应用 \hyperlink{sga6.e02.prop.1-1}{(1.1)}, 其中 $\mathcal C$ 为纤维范畴 $U\mapsto\Qcoh(U)$, $\mathcal C_0$ 为纤维范畴 $U\mapsto L(U)$. 由于 $L$ 的每个对象都伪凝聚, 纤维范畴 $\mathrm D(\Qcoh)$ 中的 $n$-伪凝聚性和伪凝聚性, 与关于 $L$ 的相应概念一致（\hyperlink{sga6.e01.cor.2-14}{I 2.14}）. 由 \hyperlink{sga6.e02.prop.2-2-3}{(2.2.3 (ii))}, \hyperlink{sga6.e02.prop.1-1}{(1.1)} 的条件 (ii) 成立, 因而（\hyperlink{sga6.e02.prop.1-1}{1.1} a)）给出 a) 中关于 $L(S)$ 的结论. 对 $\Loclib(S)$ 的断言可类似证明. b) 立即由（\hyperlink{sga6.e02.prop.1-1}{1.1} b)、c)）得出, 只须注意: 若 $E\in\Ob\mathrm D(\Qcoh(S))$ 是 $n$-伪凝聚的, 则因 $S$ 拟紧, $E$ 属于 $\mathrm D^{-}(\Qcoh(S))$. 纤维范畴 $\Qcoh$ 的对象局部属于 $L$, 当且仅当它局部自由且有限型; 因此 $L$ 在 $\Qcoh$ 中局部可提升, 且对满射的核局部稳定. 于是 c) 可由 \hyperlink{sga6.e02.prop.1-2}{(1.2)} 得出, 其中取 $\mathcal C=\Qcoh$, $\mathcal C_0=L$, $\mathcal C_1=\Loclib(S)$.

\sgaStableTarget{sga6.e02.prop.2-2-9}\textbf{命题 2.2.9.} 假设与记号同 \hyperlink{sga6.e02.thm.2-2-8}{(2.2.8)}, 再假设 $S$ 分离或诺特.
\begin{enumerate}[label=\alph*)]
\item 典范函子
\[
\mathrm K^{-,b}(L(S))/\mathrm K^{-,\delta}(L(S))\longrightarrow \mathrm D(S),
\qquad
\mathrm K^{-,b}(\Loclib(S))/\mathrm K^{-,\delta}(\Loclib(S))\longrightarrow \mathrm D(S)
\]
都是全忠实的, 且其本质像为 $\mathrm D^{b}(S)_{\coh}$（\hyperlink{sga6.e01.ex.2-15}{I 2.15}）. 此外, 若 $S$ 在 \hyperlink{sga6.e02.prop.3-7}{(3.7)} 的意义下具有有限上同调维数\footnote{第 3 节独立于本讲的其余部分.}, 例如 $S$ 的底空间诺特且 Krull 维数有限, 则函子
\[
\mathrm K^{-}(L(S))/\mathrm K^{-,\delta}(L(S))\longrightarrow
\mathrm K^{-}(\Loclib(S))/\mathrm K^{-,\delta}(\Loclib(S))\longrightarrow
\mathrm D^{-}(S)_{\coh}
\]
都是范畴等价.
\item 典范函子
\[
\mathrm K^b(\Loclib(S))/\mathrm K^{b,\delta}(\Loclib(S))\longrightarrow \mathrm D(S)
\]
是全忠实的, 且其本质像为 $\mathrm D(S)_{\parf}$（\hyperlink{sga6.e01.notation.4-9}{I 4.9}）. $E\in\Ob\mathrm D(S)$ 的完美幅包含在 $[m,n]$ 中（\hyperlink{sga6.e01.ex.4-8}{I 4.8}）的充要条件是: $E$ 在 $\mathrm D(S)$ 中同构于复形 $F$, 满足当 $i\notin[m,n]$ 时 $F^i=0$, 当 $i>m$ 时 $F^i\in\Ob L(S)$, 且 $F^m$ 局部自由且有限型.
\end{enumerate}

\emph{证明.} 由 \hyperlink{sga6.e02.prop.3-5}{(3.5)} 和 \hyperlink{sga6.e02.prop.3-7}{(3.7)}, 这是 \hyperlink{sga6.e02.thm.2-2-8}{(2.2.8)} 的直接推论. 唯一需要检验的点, 而它也平凡地由 \hyperlink{sga6.e02.prop.3-5}{(3.5)}, 分别由 \hyperlink{sga6.e02.prop.3-7}{(3.7)}, 得出, 就是: 对 $\mathrm D^{b}(\Qcoh(S))$ 的对象 $E$, 分别在 $S$ 具有 \hyperlink{sga6.e02.prop.3-7}{(3.7)} 意义下的有限上同调维数时对 $\mathrm D(\Qcoh(S))$ 的对象 $E$, 将其视为 $\mathrm D(S)$ 的对象后, \hyperlink{sga6.e02.notation.2-2-1}{(2.2.1)} 中的伪凝聚性和完美性与通常的概念一致（\hyperlink{sga6.e01.ex.2-15}{I 2.15} 及 \hyperlink{sga6.e01.notation.4-9}{4.9}）.

\sgaStableTarget{sga6.e02.rem.2-2-10}\textbf{注 2.2.10.} a) 若 $S=\operatorname{Spec}A$ 仿射, 由单个层 $\mathcal O_S$ 组成的族丰沛, 且此时 $L(S)=\Lib(S)$. 显然范畴 $\mathrm K^{-,\delta}(\Lib(S))$ 与 $\mathrm K^{-,\delta}(\Loclib(S))$ 都仅由零对象组成. 由 \hyperlink{sga6.e02.prop.2-2-9}{(2.2.9)} 可知, 若以 $\Mod(A)$ 表示 $A$-模的范畴, 并令 $\mathrm D(A)=\mathrm D(\Mod(A))$, 则函子 $R\Gamma_S:\mathrm D^{+}(S)\to\mathrm D^{+}(A)$ 诱导全忠实函子
\[
\mathrm D^{b}(S)_{\coh}\longrightarrow \mathrm D(A),
\qquad
\mathrm D(S)_{\parf}\longrightarrow \mathrm D(A),
\]
其本质像分别为三角子范畴 $\mathrm D^{b}(A)_{\coh}$ 与 $\mathrm D(A)_{\parf}$, 即上同调有界, 且关于有限型自由模的子范畴 $\Lib(A)$ 伪凝聚, 分别关于有限型投射模的子范畴 $\Loclib(A)$ 完美, 的复形所组成的范畴. 有如下范畴等价的交换图
\[
\begin{tikzcd}
\mathrm K^{-,b}(\Lib(S)) \arrow[r] \arrow[d,"\Gamma_S"'] &
\mathrm K^{-,b}(\Loclib(S)) \arrow[r] \arrow[d,"\Gamma_S"'] &
\mathrm D^b(S)_{\coh} \arrow[d,"R\Gamma_S"] \\
\mathrm K^{-,b}(\Lib(A)) \arrow[r] &
\mathrm K^{-,b}(\Loclib(A)) \arrow[r] &
\mathrm D^b(A)_{\coh},
\end{tikzcd}
\]
\[
\begin{tikzcd}
\mathrm K^b(\Loclib(S)) \arrow[r] \arrow[d,"\Gamma_S"'] &
\mathrm D(S)_{\parf} \arrow[d,"R\Gamma_S"] \\
\mathrm K^b(\Loclib(A)) \arrow[r] &
\mathrm D(A)_{\parf}.
\end{tikzcd}
\]
在上图中, 当 $S$ 具有 \hyperlink{sga6.e02.prop.3-7}{(3.7)} 意义下的有限上同调维数时, 可将 $b$ 换成 $-$.

b) 对任意概形 $S$, 由 a) 可知: 伪凝聚且上同调有界的 $\mathcal O_S$-模复形在 $\mathrm D(S)$ 中局部同构于严格伪凝聚复形, 这一点先验上远非显然（参见 \hyperlink{sga6.e01.rem.2-3-1}{I 2.3.1}）.

\sgaStableTarget{sga6.e02.sec.2.3}\subsection*{2.3. 软层}

\sgaStableTarget{sga6.e02.para.2-3-1}\textbf{2.3.1.} 回忆 [T.F.] 中的若干定义和结果. 设 $X$ 为仿紧空间. 若对 $X$ 的每个闭子集 $Y$, $F$ 在 $Y$ 上的每个截面均可延拓为 $F$ 在 $X$ 上的截面, 则称 $X$ 上的集合层 $F$ 为\emph{软层（soft sheaf）}. 阿贝尔软层无高阶上同调: 若 $F$ 软, 则对 $i>0$ 有 $H^i(X,F)=0$（T.F. II 4.4.3）.

设 $A$ 是 $X$ 上的环层. 若 $F$ 是软 $A$-模, 则 $F$ 由其截面生成（参见 \hyperlink{sga6.e02.lem.2-1-1}{2.1.1}）: 事实上, 由 $I=H^0(X,F)$ 定义的同态 $A_X^{(I)}\to F$ 是满射. 此外, 若 $A$ 软, 则每个 $A$-模都软（T.F. II 3.7.1）. 软环层的例子包括仿紧空间上的连续实值或复值函数层, 以及 $C^\infty$ 流形上的 $C^\infty$ 函数层.

\sgaStableTarget{sga6.e02.prop.2-3-2}\textbf{命题 2.3.2.} 设 $S$ 为紧空间, 配备一个软环层 $A$. 以 $\Mod(S)$ 表示左 $A$-模的范畴, 并令 $\mathrm D(S)=\mathrm D(\Mod(S))$.
\begin{enumerate}[label=\alph*)]
\item 范畴 $\Lib(S)$, 分别为 $\Loclib(S)$ \hyperlink{sga6.e02.sec.2.0}{(2.0)}, 在 $\Mod(S)$ 中可提升, 且对满射的核拟稳定, 分别稳定（\hyperlink{sga6.e01.sec.1.2}{I 1.2}）. 更一般地, 对仿紧的 $S$ 此断言也成立.
\item 设 $E\in\Ob\mathrm D(S)$ 且 $n\in\mathbb Z$. $E$ 为 $n$-伪凝聚的, 分别为伪凝聚的, 的充要条件是它关于 $\Lib(S)$ 或 $\Loclib(S)$ 具有相应性质. 典范函子
\[
\mathrm K^{-}(\Lib(S))\longrightarrow \mathrm D(S),
\qquad
\mathrm K^{-}(\Loclib(S))\longrightarrow \mathrm D(S)
\]
都是全忠实的, 且其本质像为 $\mathrm D^{-}(S)_{\coh}$（\hyperlink{sga6.e01.ex.2-15}{I 2.15}）.
\item 设 $E\in\Ob\mathrm D(S)$, $[m,n]$ 为 $\mathbb Z$ 的区间. $E$ 的完美幅包含在 $[m,n]$ 中（\hyperlink{sga6.e01.ex.4-8}{I 4.8}）的充要条件是存在拟同构 $F\to E$, 其中复形 $F$ 满足: 当 $i\notin[m,n]$ 时 $F^i=0$, 当 $i>m$ 时 $F^i\in\Ob\Lib(S)$, 且 $F^m\in\Ob\Loclib(S)$. 典范函子
\[
\mathrm K^b(\Loclib(S))\longrightarrow \mathrm D(S)
\]
是全忠实的, 且其本质像为 $\mathrm D(S)_{\parf}$（\hyperlink{sga6.e01.notation.4-9}{I 4.9}）.
\item 以 $\Mod(A(S))$ 表示左 $A(S)$-模的范畴, 以 $\mathrm D(A(S))$ 表示相应的导出范畴. 函子
\[
H^0(S,-):\Mod(S)\longrightarrow\Mod(A(S))
\]
是正合的, 且为范畴等价. 此外, 它定义范畴等价
\[
\mathrm D(S)\xrightarrow{\sim}\mathrm D(A(S)),\qquad
\mathrm D^{-}(S)_{\coh}\xrightarrow{\sim}\mathrm D^{-}(A(S))_{\coh},\qquad
\mathrm D(S)_{\parf}\xrightarrow{\sim}\mathrm D(A(S))_{\parf},
\]
这里 $\mathrm D(A(S))$ 中的伪凝聚性, 分别为完美性, 是相对于有限型自由模, 分别有限型投射模的范畴而言的. 它还给出 $\mathrm D(S)$ 中由 $n$-伪凝聚对象, 分别由完美幅包含在 $[m,n]$ 中的对象组成的满子范畴, 到 $\mathrm D(A(S))$ 的相应满子范畴的等价.
\end{enumerate}

\emph{证明.} a) 的第一条断言由（\hyperlink{sga6.e02.lem.2-1-2}{2.1.2} a)）得出: 利用 \hyperlink{sga6.e02.para.2-3-1}{(2.3.1)}, 取 $S$ 为 $S$ 上集合层的拓扑斯, $\mathcal C'_S=\Mod(S)$. 第二条由（\hyperlink{sga6.e02.prop.1-1}{1.1} a)）, 分别由（\hyperlink{sga6.e02.prop.1-2}{1.2} a)）, 得出, 其中 $\mathcal C=\Mod$, $\mathcal C_0=\Lib$, 分别为 $\Loclib$; \hyperlink{sga6.e02.prop.1-1}{(1.1)} 的条件 (ii) 由 \hyperlink{sga6.e02.para.2-3-1}{(2.3.1)} 及（\hyperlink{sga6.e02.lem.2-1-1}{2.1.1} b)）得到验证. b) 由（\hyperlink{sga6.e02.prop.1-1}{1.1} b)、c)）得出, 取相同的 $\mathcal C$ 与 $\mathcal C_0$. c) 由（\hyperlink{sga6.e02.prop.1-2}{1.2} b)、c)、d)）得出, 取 $\mathcal C=\Mod$, $\mathcal C_0=\Lib$. d) 的第一条断言利用 \hyperlink{sga6.e02.para.2-3-1}{(2.3.1)}, 由（\hyperlink{sga6.e02.lem.2-1-2}{2.1.2} b) (ii)）得出, 其中取 $\mathcal C'_S=\mathcal C_S=\Mod(S)$. d) 的其余断言由此结合 b)、c), 以及如下事实得出: $\Gamma_S$ 将有限型自由模, 分别为有限型自由模的直和因子, 变成有限型自由模, 分别有限型投射模.

\sgaStableTarget{sga6.e02.sec.2.4}\subsection*{2.4. Stein 空间}

\sgaStableTarget{sga6.e02.def.2-4-1}\textbf{定义 2.4.1.} 设 $S$ 是局部赋环拓扑斯（\hyperlink{sga6.e01.lem.2-15-1}{I 2.15.1 (ii)}）, 其终对象也记为 $S$. 若结构层 $\mathcal O_S$ 凝聚（\hyperlink{sga6.e01.def.3-1}{I 3.1}）, 且下述性质成立, 则称 $S$ 为 \emph{Stein 的}:

\emph{定理 A:} 每个凝聚 $\mathcal O_S$-模（\hyperlink{sga6.e01.def.3-1}{I 3.1}）都由其截面生成;

\emph{定理 B:} 对每个凝聚 $\mathcal O_S$-模 $F$, 当 $i>0$ 时有 $H^i(S,F)=0$.

\sgaStableTarget{sga6.e02.rem.2-4-1-1}\textbf{注 2.4.1.1.} 若 $S$ 有足够多的闭点 $s$, 使得由 $s$ 定义的理想 $\mathfrak m_s$ 凝聚, 则对每个凝聚层 $F$ 成立的关系 $H^1(S,F)=0$ 蕴含定理 A. 事实上, 设 $F$ 为凝聚 $\mathcal O_S$-模, $s$ 是 $S$ 的闭点. 有正合列
\[
0\longrightarrow \mathfrak m_sF\longrightarrow F\longrightarrow F/\mathfrak m_sF\longrightarrow0,
\]
其中 $F/\mathfrak m_sF$ 是集中于 $s$ 的凝聚层. 应用 $H^0(S,-)$ 并写出上同调正合列, 得到
\[
H^0(S,F)\longrightarrow H^0(S,F/\mathfrak m_sF)
\]
为满射. 由于
\[
H^0(S,F/\mathfrak m_sF)=F_s/\mathfrak m_sF_s,
\]
其中 $\mathfrak m_s$ 是 $\mathcal O_s$ 的极大理想, Nakayama 引理给出结论. 编写者不知道: 若仅假设 $\mathcal O_S$ 是凝聚的局部环层, 上述蕴含在一般情形下是否成立. 编写者也不知道: 对每个凝聚层 $F$ 成立的关系 $H^1(S,F)=0$ 是否蕴含定理 B.

\sgaStableTarget{sga6.e02.ex.2-4-2}\textbf{例 2.4.2.} 复解析 Stein 空间、配备由全纯函数层诱导的层的 $\mathbb C^n$ 中的紧凸子集、诺特环的谱, 以及仿射刚性解析空间, 都定义 Stein 拓扑斯.

\sgaStableTarget{sga6.e02.lem.2-4-3}\textbf{引理 2.4.3.} 设 $S$ 为 Stein 拓扑斯 \hyperlink{sga6.e02.def.2-4-1}{(2.4.1)}. 假设终对象 $S$ 拟紧 \hyperlink{sga6.e02.lem.2-1-1}{(2.1.1)}, 以 $\Coh(S)$ 表示凝聚 $\mathcal O_S$-模, 也就是有限表现模所成的阿贝尔范畴. 则:
\begin{enumerate}[label=\alph*)]
\item $\Lib(S)$, 分别为 $\Loclib(S)$, 在 $\Coh(S)$ 中可提升, 且对满射的核拟稳定, 分别稳定（\hyperlink{sga6.e01.sec.1.2}{I 1.2}）.
\item 典范函子
\[
\mathrm K^{-}(\Lib(S))\longrightarrow \mathrm K^{-}(\Loclib(S))\longrightarrow \mathrm D^{-}(\Coh(S))
\]
都是范畴等价.
\item 令 $A=\Gamma(S,\mathcal O_S)$, 以 $\Mod(A)$ 表示 $A$-模的范畴, 以 $\mathrm D(A)$ 表示相应的导出范畴. 函子
\[
\Gamma_S:\Coh(S)\longrightarrow\Mod(A)
\]
正合且全忠实, 其本质像为 $\Mod(A)$ 中由有限表现 $A$-模组成的严格满加性子范畴. 它诱导一个全忠实函子
\[
\mathrm D^{-}(\Coh(S))\longrightarrow\mathrm D(A),
\]
其本质像为由关于有限型自由模子范畴伪凝聚的复形组成的满子范畴 $\mathrm D^{-}(A)_{\coh}$.
\end{enumerate}
这里“足够多的闭点”是指足够多的形如 $i^*:S\to s$ 的纤维函子（SGA 4 IV 6.4.1）, 其中 $i_*$ 为 $S$ 的一个闭子拓扑斯的包含函子; $\mathfrak m_s$ 表示 $\mathcal O_S$ 中在 $s$ 处的纤维属于 $\mathcal O_s$ 的极大理想的截面.

\emph{证明.} 对 a) 的第一条断言, 取 $\mathcal C'_S=\Coh(S)$ 应用（\hyperlink{sga6.e02.lem.2-1-2}{2.1.2} a)）; 由定理 B, \hyperlink{sga6.e02.lem.2-1-2}{(2.1.2)} 的假设满足, 因为当 $E$ 凝聚而 $G$ 局部自由且有限型时, $\operatorname{Hom}(G,E)$ 凝聚. 因此 $\Loclib(S)$ 在 $\Coh(S)$ 中可提升, $\Lib(S)$ 更是如此. a) 的第二条断言由（\hyperlink{sga6.e02.prop.1-1}{1.1} a)）, 分别由（\hyperlink{sga6.e02.prop.1-2}{1.2} a)）, 得出, 其中 $\mathcal C_0=\Lib$, 分别为 $\Loclib$; \hyperlink{sga6.e02.prop.1-1}{(1.1)} 的条件 (ii) 由定理 A 及（\hyperlink{sga6.e02.lem.2-1-1}{2.1.1} b)）验证. b) 由（\hyperlink{sga6.e02.prop.1-1}{1.1} c)）得出, 依次取 $\mathcal C_0=\Lib$ 及 $\mathcal C_0=\Loclib$. c) 的第一条断言由（\hyperlink{sga6.e02.lem.2-1-2}{2.1.2} b)(i)）得出, 其中取 $\mathcal C'_S=\Coh(S)$. 为证明 $\Gamma_S:\mathrm D^{-}(\Coh(S))\to\mathrm D(A)$ 全忠实, 由 b), 只须证明复合函子
\[
\mathrm K^{-}(\Lib(S))\longrightarrow\mathrm D^{-}(\Coh(S))\longrightarrow\mathrm D(A)
\]
全忠实. 但此复合可分解为
\[
\mathrm K^{-}(\Lib(S))\xrightarrow{\Gamma_S}\mathrm K^{-}(\Lib(A))\longrightarrow\mathrm D(A),
\]
其中 $\Lib(A)$ 表示有限型自由 $A$-模的范畴; 第一个箭头是范畴等价, 而第二个由（\hyperlink{sga6.e01.prop.2-7}{I 2.7} b)）全忠实: 取 $\mathcal C=\Mod(A)$, $\mathcal C_0=\Lib(A)$, 并注意 $\mathrm K^{-,\delta}(\Lib(A))=0$. 这就证明了 c) 的最后一条断言.

\sgaStableTarget{sga6.e02.lem.2-4-4}\textbf{引理 2.4.4.} 在 \hyperlink{sga6.e02.lem.2-4-3}{(2.4.3)} 的假设下, 典范函子
\[
\mathrm D^b(\Coh(S))\longrightarrow\mathrm D(S)
\]
是全忠实的, 且其本质像为 $\mathrm D^b(S)_{\coh}$（\hyperlink{sga6.e01.ex.2-15}{I 2.15}）.

\emph{证明.} 应用 \hyperlink{sga6.e02.lem.2-1-3}{(2.1.3)}, 其中取 $\mathcal A=\Coh(S)$, $\mathcal B=\Mod(S)$, $\varphi:\mathcal A\to\mathcal B$ 为包含函子. 此函子全忠实, 因此由（\hyperlink{sga6.e02.lem.2-1-3}{2.1.3} a)(ii)）, 只须验证对固定的 $M\in\Ob\Coh(S)$ 及 $n>0$, 函子
\[
L\longmapsto \Ext^n_{\Mod(S)}(L,M)
\]
关于 $L\in\Ob\Coh(S)$ 余可消去. 由定理 A 及（\hyperlink{sga6.e02.lem.2-1-1}{2.1.1} b)）, 每个凝聚 $\mathcal O_S$-模都是有限型自由模的商; 当 $L$ 自由且有限型时, 由定理 B 有 $\Ext^n_{\Mod(S)}(L,M)=0$. 断言由此得出. 因而函子 $\mathrm D^b(\Coh(S))\to\mathrm D(S)$ 全忠实, 其本质像为 $\mathrm D(S)$ 中由凝聚且有界的上同调复形组成的满子范畴; 因为 $\mathcal O_S$ 凝聚, 这正是 $\mathrm D^b(S)_{\coh}$（\hyperlink{sga6.e01.cor.3-5}{I 3.5}）.

\sgaStableTarget{sga6.e02.prop.2-4-5}\textbf{命题 2.4.5.} 在 \hyperlink{sga6.e02.lem.2-4-3}{(2.4.3)} 的假设和记号下, 有:
\begin{enumerate}[label=\alph*)]
\item 典范函子
\[
\mathrm K^{-,b}(\Lib(S))\longrightarrow\mathrm K^{-,b}(\Loclib(S))\longrightarrow \mathrm D^b(S)_{\coh}
\]
都是范畴等价. 此外, 函子 $R\Gamma_S:\mathrm D^b(S)\to\mathrm D^b(A)$ 诱导全忠实函子
\[
R\Gamma_S:\mathrm D^b(S)_{\coh}\longrightarrow \mathrm D^b(A),
\]
其本质像为 $\mathrm D^b(A)_{\coh}$, 即 $\mathrm D(A)$ 中由上同调有界, 且关于有限型自由模的范畴伪凝聚的复形组成的三角子范畴.
\item 设 $E\in\Ob\mathrm D(S)$, $[m,n]$ 为 $\mathbb Z$ 的区间. 若 $E$ 的完美幅包含在 $[m,n]$ 中（\hyperlink{sga6.e01.notation.4-5}{I 4.5}）, 则 $E$ 同构于一个复形 $F$, 满足当 $i\notin[m,n]$ 时 $F^i=0$, 当 $i>m$ 时 $F^i\in\Ob\Lib(S)$, 且 $F^m\in\Ob\Loclib(S)$.
\item 典范函子
\[
\mathrm K^b(\Loclib(S))\longrightarrow \mathrm D(S)_{\parf}
\]
是范畴等价（\hyperlink{sga6.e01.notation.4-9}{I 4.9}）. 此外, $R\Gamma_S:\mathrm D^b(S)\to\mathrm D^b(A)$ 诱导全忠实函子
\[
R\Gamma_S:\mathrm D(S)_{\parf}\longrightarrow\mathrm D^b(A),
\]
其本质像为 $\mathrm D^b(A)_{\parf}$, 即 $\mathrm D(A)$ 中由关于有限型投射模范畴完美的复形组成的三角子范畴.
\end{enumerate}

\emph{证明.} a) 由（\hyperlink{sga6.e02.lem.2-4-3}{2.4.3} b)、c)）, 函子
\[
\mathrm K^{-,b}(\Lib(S))\longrightarrow\mathrm K^{-,b}(\Loclib(S))\longrightarrow
\mathrm D^b(\Coh(S))\longrightarrow\mathrm D^b(A)_{\coh}
\]
都是范畴等价. 由于全忠实函子 $\Gamma_S:\mathrm D^b(\Coh(S))\to\mathrm D(A)$ 分解为
\[
\mathrm D^b(\Coh(S))\longrightarrow\mathrm D^b(S)_{\coh}\xrightarrow{R\Gamma_S}\mathrm D(A),
\]
结论便由 \hyperlink{sga6.e02.lem.2-4-4}{(2.4.4)} 得出.

b) 假设 $\operatorname{parf.amp}(E)\subset[m,n]$. 特别地, $E$ 伪凝聚; 由 a), 以同构对象代替 $E$ 后, 可以假设当 $i>m$ 时 $E^i\in\Ob\Lib(S)$, 且当 $i>n$ 时 $E^i=0$. 再作截断, 还可假设当 $i<m$ 时 $E^i=0$. 取 $\mathcal C=\Mod$, $\mathcal C_0=\Lib$ 应用（\hyperlink{sga6.e01.lem.4-16}{I 4.16}）, 便有 $E^m\in\Ob\Loclib(S)$.

c) 由（\hyperlink{sga6.e02.lem.2-4-3}{2.4.3} a)）, 以及（\hyperlink{sga6.e02.lem.2-1-1}{2.1.1} b)）结合定理 A, 对 $\mathcal C=\Coh(S)$、$\mathcal C_0=\Lib$, 假设 \hyperlink{sga6.e02.prop.1-1}{(1.1 (i))} 以及相应的 \hyperlink{sga6.e02.prop.1-1}{(1.1 (ii))} 均得到验证. 因而（\hyperlink{sga6.e02.prop.1-2}{1.2} c)、d)）表明函子
\[
\mathrm K^b(\Loclib(S))\longrightarrow\mathrm D^b(\Coh(S))
\]
全忠实. 因此由 \hyperlink{sga6.e02.lem.2-4-4}{(2.4.4)}, 函子
\[
\mathrm K^b(\Loclib(S))\longrightarrow\mathrm D(S),
\]
也全忠实, 且由 b), 其本质像为 $\mathrm D(S)_{\parf}$. 采用类似于（\hyperlink{sga6.e02.lem.2-4-3}{2.4.3} c)）证明中的论证, 可知由 $\Gamma_S$ 诱导的函子 $\mathrm K^b(\Loclib(S))\to\mathrm D(A)$ 全忠实, 其本质像为 $\mathrm D^b(A)_{\parf}$; 于是如 a) 中一样完成证明.

\sgaStableTarget{sga6.e02.rem.2-4-6}\textbf{注 2.4.6.} 编写者不知道在 \hyperlink{sga6.e02.lem.2-4-4}{(2.4.4)} 中是否可以用 $\mathrm D^{-}$ 代替 $\mathrm D^b$, 因而也不知道在用 $\mathrm K^{-}$ 代替 $\mathrm K^{-,b}$、用 $\mathrm D^{-}$ 代替 $\mathrm D^b$ 后, （\hyperlink{sga6.e02.prop.2-4-5}{2.4.5} a)）是否仍成立.

最后给出 \hyperlink{sga6.e02.prop.2-4-5}{(2.4.5)} 的一个有趣补充. 先作一个定义.

\sgaStableTarget{sga6.e02.def.2-4-7}\textbf{定义 2.4.7.} 所谓\emph{复伪解析空间（complex pseudo-analytic space）}, 是指局部同构于形如 $(Y,\mathcal O_X|Y)$ 的赋环空间的赋环空间, 其中 $X$ 为复解析空间, $Y$ 为 $X$ 的闭子集.

因此复伪解析空间的结构层凝聚. 若相应的拓扑斯是 Stein 的 \hyperlink{sga6.e02.def.2-4-1}{(2.4.1)}, 则称该复伪解析空间为 Stein 的.



% 第 II 讲续篇的宏
\providecommand{\Qcoh}{\operatorname{Qcoh}}
\providecommand{\Mod}{\operatorname{Mod}}
\providecommand{\Coh}{\operatorname{Coh}}
\providecommand{\Ob}{\operatorname{Ob}}
\providecommand{\Hom}{\operatorname{Hom}}
\providecommand{\Ext}{\operatorname{Ext}}
\providecommand{\Spec}{\operatorname{Spec}}
\providecommand{\qcoh}{\mathrm{qcoh}}
\providecommand{\parf}{\mathrm{parf}}
\providecommand{\coh}{\mathrm{coh}}
\providecommand{\Loclib}{\operatorname{Loclib}}
\providecommand{\Lib}{\operatorname{Lib}}
\providecommand{\im}{\operatorname{Im}}

\sgaStableTarget{sga6.e02.expose.occ2}\section*{第 II 讲. 整体消解（续）}

\sgaStableTarget{sga6.e02.prop.2-4-8}\textbf{命题 2.4.8.} 设 $S$ 是紧的 Stein 复伪解析空间 \hyperlink{sga6.e02.def.2-4-7}{(2.4.7)}. 假设对 $S$ 的每个开集 $U$ 及 $U$ 的每个约化闭子空间 $X$, $X$ 的正规化局部连通. 例如, 当 $S$ 在空间 $\mathbb R^n$ 中局部由实解析方程或不等式定义时, 此条件满足. 则
\[
 A=\Gamma(S,\mathcal O_S)
\]
是诺特环.

\emph{证明.} 只须证明: 若 $M$ 是有限表现 $A$-模, 则 $M$ 的所有子模, 或等价地 $M$ 的所有有限型子模, 满足极大条件. 事实上, 若 $N$ 是 $M$ 的有限型子模, 则由 \hyperlink{sga6.e01.schol.2-6}{I 2.6} 及 \hyperlink{sga6.e01.prop.2-9}{2.9}, $M/N$ 有限表现, 从而由等价（\hyperlink{sga6.e02.lem.2-4-3}{2.4.3} c)）, $N$ 本身也有限表现. 再用凝聚 $\mathcal O_S$-模的范畴与有限表现 $A$-模的范畴之间的等价（\hyperlink{sga6.e02.lem.2-4-3}{2.4.3} c)）, 并利用 $S$ 的紧性, 问题便归结为证明下述引理.

\sgaStableTarget{sga6.e02.lem.2-4-8-1}\textbf{引理 2.4.8.1.} 设 $S$ 是满足 \hyperlink{sga6.e02.prop.2-4-8}{(2.4.8)} 中局部连通性条件的复伪解析空间. 设 $F$ 为凝聚 $\mathcal O_S$-模, $(E_n)$ 为 $F$ 的凝聚子 $\mathcal O_S$-模的递增序列, $s$ 为 $S$ 的点. 则存在 $s$ 的开邻域 $U$, 使序列 $(E_n|U)$ 最终稳定.

\emph{证明.} 由于环 $\mathcal O_{S,s}$ 诺特, 存在整数 $n_0$, 使得当 $n\geq n_0$ 时 $(E_n)_s=(E_{n_0})_s$. 以序列 $(E_n/E_{n_0})$ 代替 $(E_n)$ 后, 可以假设对每个 $n$ 有 $(E_n)_s=0$. 另一方面, $F_s=M$ 是有限 $\mathcal O_{S,s}$-模; 因而它有有限滤过
\[
 M=M_0\supset M_1\supset\cdots\supset M_q=0
\]
满足
\[
 M_i/M_{i+1}\simeq \mathcal O_s/\mathfrak p_i,
\]
其中 $\mathfrak p_i$ 为 $\mathcal O_s$ 的素理想. 由于有限型 $\mathcal O_s$-模的范畴等价于凝聚 $\mathcal O_S$-模在 $s$ 处的芽的范畴, 缩小 $s$ 的邻域后, 可假设 $M_i$, 分别为 $\mathfrak p_i$, 是 $F$ 的凝聚子 $\mathcal O_S$-模 $F_i$, 分别为 $\mathcal O_S$ 的凝聚理想 $\mathfrak p_i$, 在 $s$ 处的茎, 且有
\[
 F=F_0\supset F_1\supset\cdots\supset F_q=0,
 \qquad F_i/F_{i+1}\simeq\mathcal O_S/\mathfrak p_i.
\]
因此, 证明序列 $(E_n)$ 在 $s$ 的某邻域中最终稳定, 等价于对这一滤过的各个分次商上诱导的序列证明同一断言. 因而可以假设 $F=\mathcal O_S$, 且 $\mathcal O_{S,s}$ 为整环; 此时 $E_n$ 构成 $\mathcal O_S$ 的理想的递增序列, 满足 $(E_n)_s=0$.

由于 $\mathcal O_{S,s}$ 为整环, 特别地它约化, 由 Oka 定理, 在 $s$ 处局部化后还可以假设 $S$ 约化. 需要证明存在 $s$ 的开邻域 $U$, 使得当 $n$ 充分大时 $E_n|U=0$. 为此考虑 $S$ 的正规化 $\widetilde S$. 它是有限 $\mathcal O_S$-代数 $\widetilde{\mathcal O}_S$ 的解析谱, 该代数在每个 $t\in S$ 处的茎为 $\mathcal O_{S,t}$ 的整闭包. 由于 $E_n\subset E_n\widetilde{\mathcal O}_S$, 只须证明 $n$ 充分大时 $E_n\widetilde{\mathcal O}_S=0$. 因而可归结为 $S$ 正规的情形.

此时, 由于 $(E_n)_s=0$, 由 $E_n$ 定义的闭浸入
\[
 i_n:S_n\longrightarrow S
\]
在 $s$ 处平展. 但由于 $S$ 正规, 由下述引理, $i_n$ 平展的点集既开又闭. 因此其在 $S$ 中的像同时为开集和闭集. 从而 $E_n$ 在 $s$ 所属的 $S$ 的连通分支上的限制为零; 由于 $S$ 局部连通, 此分支为开集. 引理得证.

\sgaStableTarget{sga6.e02.lem.2-4-8-2}\textbf{引理 2.4.8.2}（参见 EGA IV 18.10.3）. 设 $Y$ 是正规伪解析空间\footnote{只须 $Y$ 的局部环为整环.}, $f:X\to Y$ 为非分歧态射. $X$ 中 $f$ 平展的点所成的集合既开又闭.

\emph{证明.} $f$ 平展的点集显然为开集, 因为 $f$ 在某点平展, 就是说它在该点的某邻域中为同构（[8] VI）. 只须证明 $f$ 不平展的点集为开集. 由于非分歧态射局部为浸入（[8] VI）, 可以假设 $f$ 为由凝聚理想 $I$ 定义的闭浸入. 设 $x$ 是 $X$ 的一个点, 且 $I_x\ne0$. 取包含 $I_x$ 的 $\mathcal O_{Y,x}$ 的素理想 $\mathfrak p$. 有
\[
 \dim(\mathcal O_{Y,x})_{\mathfrak p}>0,
\]
否则因 $\mathcal O_{Y,x}$ 为整环, 必有 $\mathfrak p=0$. 因而 $Y$ 在 $x$ 处的余维大于 $0$, 从而在 $x$ 的某邻域中余维也大于 $0$（[10] III A, 引理 3 的推论 1）; 断言由此得出.

\sgaStableTarget{sga6.e02.sec.3}\section*{3. 概形上拟凝聚层的补充}\footnote{本节主要归功于 J.-L. Verdier.}

\sgaStableTarget{sga6.e02.para.3-0}\textbf{3.0.} 本节固定一个拟紧拟分离概形 $S$. 典范的全忠实包含函子
\[
 \iota:\Qcoh(S)\longrightarrow\Mod(S)\footnotemark
\]
延拓为由 $\mathrm D(\Qcoh(S))$ 到 $\mathrm D(S)=\mathrm D(\Mod(S))$ 的函子, 仍记为 $\iota$, 下面将研究此函子\footnote{在不会引起混淆的公式中, 将省略记号 $\iota$.}.\footnotetext{此处 $\Mod(S)$, 分别为 $\Qcoh(S)$, 表示 $\mathcal O_S$-模, 分别为拟凝聚 $\mathcal O_S$-模的范畴.}

\sgaStableTarget{sga6.e02.lem.3-1}\textbf{引理 3.1.} 范畴 $\Qcoh(S)$ 是满足 AB5 的阿贝尔范畴, 其中有限型拟凝聚层构成一个小的生成元族. 特别地（Tohoku 1.10.1）, $\Qcoh(S)$ 有足够多的内射对象.

\emph{证明.} 关于生成元的断言立即由下述事实得出: $S$ 上的每个拟凝聚层都是其有限型拟凝聚子层的滤过归纳极限（EGA I 9.4.9 及 IV 1.7.7）. 其余断言是显然的.

\sgaStableTarget{sga6.e02.rem.3-1-1}\textbf{注 3.1.1.} 编写者不知道: 若只假设 $S$ 拟分离, $\Qcoh(S)$ 是否仍有足够多的内射对象.

\sgaStableTarget{sga6.e02.lem.3-2}\textbf{引理 3.2.} 函子
\[
 \iota:\Qcoh(S)\longrightarrow\Mod(S)
\]
有右伴随函子
\[
 Q:\Mod(S)\longrightarrow\Qcoh(S),
\]
称为凝聚化函子（coherer）.

\emph{证明.} 若 $S$ 仿射, $Q$ 的存在性是直接的: 取 $Q=\widetilde{\Gamma_S}$（EGA I 1.3.7）. 现设 $f:U\to S$ 是 $S$ 的一个仿射开集. 显然, 对 $E\in\Ob\Qcoh(S)$ 和 $F\in\Ob\Mod(U)$, 有典范同构
\sgaStableTarget{sga6.e02.eq.1}\begin{equation}\tag{1}
 \Hom(E,f_*F)\simeq \Hom(E,f_*Q(F)).
\end{equation}
取 $S$ 的有限仿射开覆盖 $(f_i:U_i\to S)$; 对每对 $(i,j)$, 取 $U_i\cap U_j$ 的有限仿射开覆盖 $(U_{ijk}\to U_i\cap U_j)_k$, 这是可能的, 因为 $S$ 拟分离. 对 $F\in\Ob\Mod(S)$, 有正合列
\sgaStableTarget{sga6.e02.eq.2}\begin{equation}\tag{2}
 0\longrightarrow F\longrightarrow \prod_i f_{i*}(F|U_i)
 \longrightarrow \prod_{i,j,k}f_{ijk*}(F|U_{ijk}),
\end{equation}
其中 $f_{ijk}$ 表示 $U_{ijk}$ 到 $S$ 的包含映射. 若 $g:V\to U$ 是 $S$ 的仿射开集之间的包含映射, 则有显然的自然箭头
\[
 Q(F|U)\longrightarrow g_*Q(F|V).
\]
因而由 (2) 得到箭头
\sgaStableTarget{sga6.e02.eq.3}\begin{equation}\tag{3}
 \prod_i f_{i*}Q(F|U_i)\longrightarrow
 \prod_{i,j,k}f_{ijk*}Q(F|U_{ijk}),
\end{equation}
以 $QF$ 表示其核. 若 $E\in\Ob\Qcoh(S)$, 将函子 $\Hom(E,\cdot)$ 应用于 (2), 并使用 (1), 得到关于 $E$ 函子性的同构
\[
 \Hom(E,F)\simeq\Hom(E,QF).
\]
这就证明了 $Q$ 的存在性; 由于它为伴随函子, 其唯一性是显然的.

\sgaStableTarget{sga6.e02.lem.3-3}\textbf{引理 3.3.} \begin{enumerate}[label=\alph*)]
\item 凝聚化函子是加性函子, 且将内射对象变成内射对象.
\item 设
\[
 RQ:\mathrm D^+(S)\longrightarrow \mathrm D^+(\Qcoh(S))
\]
为凝聚化函子的右导出函子. 对 $E\in\Ob\mathrm D(\Qcoh(S))$ 和 $F\in\Ob\mathrm D^+(S)$, 有典范的函子性同构
\sgaStableTarget{sga6.e02.eq.i}\sgaStableTarget{sga6.e02.eq.ii}\begin{align}
 R\Hom_{\Mod}(E,F)&\simeq R\Hom_{\Qcoh}(E,RQ(F)), \tag{i}\\
 \Hom_{\Mod}(E,F)&\simeq \Hom_{\Qcoh}(E,RQ(F)). \tag{ii}
\end{align}
\end{enumerate}

\emph{证明.} a) 是定义的直接推论. 对 b), 只须建立同构 (i), 因为 (ii) 可由取上同调得出. 为此可假设 $F$ 的各项内射. 于是由 a), $QF$ 的各项也内射, 同构 (i) 便直接由 $Q$ 的定义得出.

\sgaStableTarget{sga6.e02.lem.3-4}\textbf{引理 3.4.} 设 $f:X\to Y$ 为拟紧拟分离概形之间的态射. 对 $E\in\Ob\Mod(X)$, 有典范的函子性同构
\[
 Q_Y f_*(E)\simeq f_*Q_X(E).
\]

\emph{证明.} 由于 $f$ 拟紧拟分离（EGA IV 1.2.3 及 1.2.4）, $f_*$ 将拟凝聚层变成拟凝聚层（EGA III 1.4.10 及 IV 1.7.21）. 因而由 $\Mod(X)$ 到 $\Qcoh(Y)$ 的函子 $Q_Yf_*$ 与 $f_*Q_X$ 都是函子
\[
 f^*: \Qcoh(Y)\longrightarrow \Mod(X),
\]
的右伴随, 因而典范同构.

\sgaStableTarget{sga6.e02.prop.3-5}\textbf{命题 3.5.} 若 $S$ 诺特, 或拟紧且分离, 则函子
\[
 \iota:\mathrm D^+(\Qcoh(S))\longrightarrow \mathrm D^+(S)
\]
是全忠实的, 且其本质像为 $\mathrm D^+(S)$ 中由上同调拟凝聚的复形组成的满子范畴 $\mathrm D^+(S)_{\qcoh}$.

\emph{证明.} 可以只在 $S$ 分离时证明此断言, 因为诺特情形已知. 显然, 这等价于证明下述两个断言.

\begin{enumerate}[label=(3.5.\arabic*)]
\item 对 $E\in\Ob\mathrm D^+(\Qcoh(S))$, 伴随箭头
\[
 E\longrightarrow RQ(E)
\]
是同构.
\item 对 $E\in\Ob\mathrm D^+(S)_{\qcoh}$, 伴随箭头
\[
 RQ(E)\longrightarrow E
\]
是同构.
\end{enumerate}

\emph{\hyperlink{sga6.e02.prop.3-5}{(3.5.1)} 的证明.} 由直接的逐步分解论证, 可归结为证明拟凝聚层对凝聚化函子无高阶上同调, 即对 $E\in\Ob\Qcoh(S)$, 当 $i>0$ 时有 $R^iQ(E)=0$, 且伴随箭头 $E\to QE$ 是同构. 若 $S$ 仿射, 这由 \hyperlink{sga6.e02.lem.3-2}{(3.2)} 证明中使用的公式 $Q=\widetilde{\Gamma_S}$ 得出. 一般地, 用有限个仿射开集 $U_i$ 覆盖 $S$. 由于 $S$ 分离, 包含映射
\[
 f_{i_0\ldots i_p}:U_{i_0}\cap\cdots\cap U_{i_p}\longrightarrow S
\]
都是仿射态射. 设 $E\in\Ob\Qcoh(S)$, 考虑由该覆盖定义的 $E$ 的消解
\sgaStableTarget{sga6.e02.eq.unnamed.occ3}\begin{equation}\tag{*}
 0\longrightarrow E\longrightarrow \prod_i f_{i*}f_i^*E\longrightarrow\cdots\longrightarrow
 \prod_{i_0\ldots i_p}f_{i_0\ldots i_p*}f_{i_0\ldots i_p}^*E\longrightarrow\cdots .
\end{equation}
它给出谱序列
\sgaStableTarget{sga6.e02.eq.unnamed.occ4}\begin{equation}\tag{**}
 E_1^{pq}=\prod_{i_0\ldots i_p}R^qQ\bigl(f_{i_0\ldots i_p*}f_{i_0\ldots i_p}^*E\bigr)
 \Longrightarrow R^{p+q}Q(E).
\end{equation}
由于 $f_{i_0\ldots i_p}$ 仿射, \hyperlink{sga6.e02.lem.3-4}{(3.4)} 对每个 $q$ 给出
\sgaStableTarget{sga6.e02.eq.unnamed.occ5}\begin{equation}\tag{***}
 R^qQ\bigl(f_{i_0\ldots i_p*}f_{i_0\ldots i_p}^*E\bigr)
 \simeq f_{i_0\ldots i_p*}R^qQ\bigl(f_{i_0\ldots i_p}^*E\bigr).
\end{equation}
特别地, 由已经处理过的仿射情形, 当 $q>0$ 时 $E_1^{pq}=0$. 由 (**) 推出同构
\[
 R^nQ(E)\simeq H^n(QE')\quad(n\geq0),
\]
其中 $E'$ 是由正合列 (*) 定义的 $E$ 的消解. 但由 (***), 当 $n\geq0$ 时有 $QE'{}^n\simeq E'{}^n$; 因而最终得到
\[
 R^nQ(E)=0\quad(n>0),
\]
以及 $Q(E)\simeq E$. 这就证明了 \hyperlink{sga6.e02.prop.3-5}{(3.5.1)}.

\emph{\hyperlink{sga6.e02.prop.3-5}{(3.5.2)} 的证明.} 谱序列
\[
 R^pQ(H^q(E))\Longrightarrow R^{p+q}Q(E)
\]
立即将问题归结为 $E\in\Ob\Qcoh(S)$ 的情形. 伴随箭头的复合
\[
 E\longrightarrow RQE\longrightarrow E
\]
是恒等映射, 于是断言由 \hyperlink{sga6.e02.prop.3-5}{(3.5.1)} 得出. \hyperlink{sga6.e02.prop.3-5}{(3.5)} 得证.

\sgaStableTarget{sga6.e02.rem.3-6}\textbf{注 3.6.} 若省略“分离或诺特”的假设, \hyperlink{sga6.e02.prop.3-5}{(3.5)} 的结论未必成立; J.-L. Verdier 给出了一个精巧的反例, 参见本讲附录 I.

\sgaStableTarget{sga6.e02.prop.3-7}\textbf{命题 3.7.} 在 \hyperlink{sga6.e02.prop.3-5}{(3.5)} 的假设下, 再假设存在整数 $N$, 使得对每个 $F\in\Ob\Mod(S)$ 及 $S$ 的每个拟紧开集 $U$, 都有
\[
 H^i(U,F)=0\quad\text{当 }i>N.
\]
则:
\begin{enumerate}[label=\alph*)]
\item 凝聚化函子的上同调维数有限.
\item 函子
\[
 \iota:\mathrm D(\Qcoh(S))\longrightarrow \mathrm D(S)
\]
是全忠实的, 且其本质像为 $\mathrm D(S)$ 中由上同调拟凝聚的复形组成的满子范畴 $\mathrm D(S)_{\qcoh}$.
\end{enumerate}

\emph{a) 的证明.} 若 $S$ 仿射, 由公式 $Q=\widetilde{\Gamma_S}$ 可知 $Q$ 的上同调维数 $\leq N$. 现转向一般情形. 证明存在整数 $a$, 使得对每个 $E\in\Ob\Mod(S)$ 及每个 $q>a$, 有 $R^qQ(E)=0$. 假设蕴含每个 $E\in\Ob\Mod(S)$ 都有长度为 $N$ 的有限右消解 $E'$, 其中 $E'{}^i$ 无高阶上同调, 即对 $S$ 的每个开集 $U$ 及每个 $q>0$, 有 $H^q(U,E'{}^i)=0$. 因而只须证明存在整数 $b$, 使得
\[
 R^qQ(E)=0
\]
对每个无高阶上同调的 $\mathcal O_S$-模 $E$ 及每个 $q>b$ 成立; 此时可取 $a=b+N$.

取由 $N'$ 个仿射开集构成的 $S$ 的有限覆盖 $(U_i)$. 设 $E$ 是无高阶上同调的 $\mathcal O_S$-模, $C'(E)$ 是由覆盖 $(U_i)$ 定义的 $E$ 的消解. 分两种情形.

若 $S$ 分离, 则 $C'(E)$ 的各项为形如 $f_*f^*E$ 的项之和, 其中 $f:U\to S$ 为仿射开集的包含映射, 因而为仿射态射. 有
\[
 RQ(Q_Xf_*)f^*E\simeq RQ\,Rf_*f^*E\simeq RQ\,f_*f^*E,
\]
因为 $f^*E$ 无高阶上同调. 另一方面, 由 \hyperlink{sga6.e02.lem.3-4}{(3.4)},
\[
 R(Q_Xf_*)f^*E\simeq R(f_*Q_X)f^*E\simeq f_*RQf^*E,
\]
因为 $f$ 仿射. 从而
\[
 R^nQ(f_*f^*E)\simeq f_*R^nQ(f^*E)\quad(n\geq0).
\]
由仿射情形,
\[
 R^nQ(f_*f^*E)=0\quad(n>N).
\]
由于 $C'(E)$ 的长度为 $N'$, 得到
\[
 R^nQ(E)=0\quad(n>N+N').
\]

若 $S$ 诺特, 则 $C'(E)$ 的各项为同样形式的项 $f_*f^*E$ 之和, 但此时 $f:U\to S$ 为拟紧且分离的开集的包含映射. 与第一种情形一样, 有
\[
 R(Q_Xf_*)f^*E\simeq RQ\,f_*f^*E.
\]
由于 $U$ 诺特, $\Qcoh(U)$ 有足够多的在 $\Mod(U)$ 中内射的对象; 因而在由 \hyperlink{sga6.e02.prop.3-5}{(3.5)} 将 $\mathrm D^+(\Qcoh(U))$ 与 $\mathrm D^+(\Qcoh(S))$ 分别视为 $\mathrm D^+(U)$ 与 $\mathrm D^+(S)$ 的满子范畴时, $f_*:\Qcoh(U)\to\Qcoh(S)$ 的右导出函子由 $Rf_*:\mathrm D^+(U)\to\mathrm D^+(S)$ 诱导. 因此
\[
 R(f_*Q_X)f^*E\simeq Rf_*RQf^*E.
\]
推出谱序列
\[
 E_2^{pq}=R^pf_*R^qQ(f^*E)\Longrightarrow R^{p+q}(f_*Q)f^*E.
\]
由第一种情形, 存在只依赖于 $U$ 的整数 $n(U)$, 使得 $Q$ 在 $\Mod(U)$ 上的上同调维数 $<n(U)$; 因而当 $q>n(U)$ 时 $E_2^{pq}=0$. 由 $S$ 的上同调维数假设, 当 $p>N$ 时也有 $E_2^{pq}=0$. 因此
\[
 R^n(f_*Q)f^*E=0\quad(n>N+N''),
\]
